\chapter{Les composés carbonylés — aldéhydes et cétones}

\textit{Les composés carbonylés, caractérisés par la présence du groupe fonctionnel carbonyle \ce{C=O}, occupent une place centrale en chimie organique. On distingue deux grandes familles : les aldéhydes, où le carbonyle est en bout de chaîne, et les cétones, où il est à l'intérieur de la chaîne carbonée. Leurs propriétés chimiques, notamment l'addition nucléophile et l'oxydation, en font des intermédiaires de synthèse essentiels. Ce chapitre présente leur nomenclature, leurs méthodes d'obtention, leurs tests d'identification et leurs réactions caractéristiques.}

\section{Structure du groupe carbonyle}

\begin{definition}[Groupe carbonyle]
Le groupe carbonyle est un groupe fonctionnel constitué d'un atome de carbone lié par une double liaison à un atome d'oxygène : \ce{C=O}. Le carbone est hybridé $sp^2$, ce qui confère une géométrie plane trigonale autour de lui.
\end{definition}

\begin{propriete}[Polarité du groupe carbonyle]
En raison de la différence d'électronégativité entre le carbone et l'oxygène, la liaison \ce{C=O} est fortement polarisée : l'oxygène porte une charge partielle négative ($\delta^-$) et le carbone une charge partielle positive ($\delta^+$). Cette polarité rend le carbone électrophile (sensible aux attaques nucléophiles) et l'oxygène basique.
\end{propriete}

\begin{illus}
\chemfig{C(=[1]O)-[2]R_1-[0]R_2}
\fig{Représentation générale du groupe carbonyle. $R_1$ et $R_2$ sont des groupes alkyles ou aryles.}
\end{illus}

\section{Aldéhydes et cétones : définitions et nomenclature}

\subsection{Définitions}

\begin{definition}[Aldéhyde]
Un aldéhyde est un composé carbonylé dans lequel le carbone du groupe carbonyle est lié à au moins un atome d'hydrogène. La formule générale est \ce{R-CHO} (ou \ce{R-CH=O}), où R est un groupe alkyle ou aryle. Le groupe fonctionnel est le groupe formyle \ce{-CHO}.
\end{definition}

\begin{definition}[Cétone]
Une cétone est un composé carbonylé dans lequel le carbone du groupe carbonyle est lié à deux atomes de carbone (groupes alkyles ou aryles). La formule générale est \ce{R-CO-R$'$}, où R et R' sont des groupes alkyles ou aryles. Le groupe fonctionnel est le groupe carbonyle \ce{C=O}.
\end{definition}

\begin{remarque}
Le formaldéhyde (méthanal) \ce{HCHO} est le seul aldéhyde où R = H. L'acétone (propanone) \ce{CH3COCH3} est la cétone la plus simple.
\end{remarque}

\subsection{Nomenclature}

\begin{loi}[Nomenclature des aldéhydes]
\begin{itemize}
    \item On repère la chaîne carbonée la plus longue contenant le groupe \ce{-CHO}.
    \item Le nom de l'alcane correspondant est modifié en remplaçant le \textit{-e} final par \textit{-al}.
    \item Le carbone du groupe \ce{-CHO} est toujours le carbone numéro 1.
    \item Les substituants sont nommés et numérotés comme d'habitude.
\end{itemize}
\end{loi}

\begin{exemple}
\begin{itemize}
    \item \ce{HCHO} : méthanal (formaldéhyde)
    \item \ce{CH3CHO} : éthanal (acétaldéhyde)
    \item \ce{CH3CH2CHO} : propanal
    \item \ce{CH3CH(CH3)CHO} : 2-méthylpropanal
    \item \chemfig{CHO-*6(-=-=-=)} : benzaldéhyde
\end{itemize}
\end{exemple}

\begin{loi}[Nomenclature des cétones]
\begin{itemize}
    \item On repère la chaîne carbonée la plus longue contenant le groupe \ce{C=O}.
    \item Le nom de l'alcane correspondant est modifié en remplaçant le \textit{-e} final par \textit{-one}.
    \item On indique la position du carbonyle par le numéro le plus petit possible.
    \item Les substituants sont nommés et numérotés.
\end{itemize}
\end{loi}

\begin{exemple}
\begin{itemize}
    \item \ce{CH3COCH3} : propanone (acétone)
    \item \ce{CH3COCH2CH3} : butan-2-one
    \item \ce{CH3CH2COCH2CH3} : pentan-3-one
    \item \chemfig{*6(-(-[:30]O)=(-[:150]CH_3)-(-[:210]CH_3)-=--)} : 4-méthylpent-3-èn-2-one (exemple de cétone insaturée)
\end{itemize}
\end{exemple}

\begin{aretenir}
Pour nommer un aldéhyde, on utilise le suffixe \textit{-al} ; pour une cétone, le suffixe \textit{-one} avec l'indice de position du carbonyle.
\end{aretenir}

\section{Obtention des aldéhydes et cétones}

\subsection{Oxydation des alcools}

\begin{experience}[Oxydation ménagée des alcools]
L'oxydation ménagée des alcools primaires et secondaires permet d'obtenir respectivement des aldéhydes et des cétones. L'oxydant utilisé est souvent le dichromate de potassium \ce{K2Cr2O7} en milieu acide ou le permanganate de potassium \ce{KMnO4} dilué à froid.
\end{experience}

\begin{propriete}[Oxydation des alcools primaires]
Un alcool primaire \ce{R-CH2OH} est oxydé en aldéhyde \ce{R-CHO}, puis, si l'oxydation se poursuit, en acide carboxylique \ce{R-COOH}.
\[
\ce{R-CH2OH + [O] -> R-CHO + H2O}
\]
\[
\ce{R-CHO + [O] -> R-COOH}
\]
\end{propriete}

\begin{propriete}[Oxydation des alcools secondaires]
Un alcool secondaire \ce{R-CHOH-R$'$} est oxydé en cétone \ce{R-CO-R$'$}. L'oxydation s'arrête à ce stade car les cétones sont résistantes à l'oxydation douce.
\[
\ce{R-CHOH-R$'$ + [O] -> R-CO-R$'$ + H2O}
\]
\end{propriete}

\begin{remarque}
Les alcools tertiaires ne s'oxydent pas dans ces conditions.
\end{remarque}

\begin{exemple}
\begin{itemize}
    \item \ce{CH3CH2OH + [O] -> CH3CHO + H2O} (éthanol $\rightarrow$ éthanal)
    \item \ce{CH3CHOHCH3 + [O] -> CH3COCH3 + H2O} (propan-2-ol $\rightarrow$ propanone)
\end{itemize}
\end{exemple}

\subsection{Autres méthodes}

\begin{propriete}[Obtention par oxydation des alcènes]
L'ozonolyse des alcènes suivie d'une réduction (par \ce{Zn/H2O} ou \ce{Me2S}) permet d'obtenir des aldéhydes ou des cétones selon la structure de l'alcène.
\end{propriete}

\begin{propriete}[Obtention par hydrolyse des dérivés halogénés]
L'hydrolyse des gem-dihalogénures (composés avec deux halogènes sur le même carbone) en milieu basique donne des aldéhydes ou des cétones.
\end{propriete}

\section{Propriétés physiques}

\begin{propriete}[Propriétés physiques]
\begin{itemize}
    \item Les aldéhydes et cétones sont polaires en raison de la liaison \ce{C=O}.
    \item Les premiers termes (méthanal, éthanal, propanone) sont miscibles à l'eau en raison de liaisons hydrogène avec l'eau.
    \item Le point d'ébullition est plus élevé que celui des alcanes correspondants, mais plus bas que celui des alcools (pas de liaison hydrogène entre molécules carbonylées).
    \item Les aldéhydes et cétones sont généralement liquides à température ambiante (sauf le méthanal qui est gazeux).
\end{itemize}
\end{propriete}

\section{Réactions caractéristiques et tests d'identification}

\subsection{Test à la 2,4-dinitrophénylhydrazine (2,4-DNPH)}

\begin{experience}[Test à la 2,4-DNPH]
La 2,4-dinitrophénylhydrazine (2,4-DNPH) réagit avec les aldéhydes et les cétones pour former une 2,4-dinitrophénylhydrazone, un précipité de couleur jaune à orange.
\end{experience}

\begin{propriete}[Réaction avec la 2,4-DNPH]
La réaction est une addition-élimination (condensation) :
\[
\ce{R-CO-R$'$ + H2N-NH-C6H3(NO2)2 -> R-C(=N-NH-C6H3(NO2)2)-R$'$ + H2O}
\]
Le précipité formé est caractéristique de la présence d'un groupe carbonyle.
\end{propriete}

\begin{remarque}
Ce test est positif pour tous les aldéhydes et toutes les cétones. Il ne permet pas de distinguer les deux familles.
\end{remarque}

\subsection{Réactif de Schiff}

\begin{experience}[Test au réactif de Schiff]
Le réactif de Schiff (fuchsine décolorée par le dioxyde de soufre) donne une coloration rose-magenta en présence d'un aldéhyde. Les cétones ne réagissent pas.
\end{experience}

\begin{propriete}[Réaction avec le réactif de Schiff]
Seuls les aldéhydes redonnent la couleur rose au réactif de Schiff. Ce test permet de distinguer un aldéhyde d'une cétone.
\end{propriete}

\subsection{Liqueur de Fehling}

\begin{experience}[Test à la liqueur de Fehling]
La liqueur de Fehling est une solution alcaline contenant des ions cuivre(II) complexes par l'acide tartrique (tartrate de cuivre). En présence d'un aldéhyde, chauffé, il se forme un précipité rouge brique d'oxyde de cuivre(I) \ce{Cu2O}.
\end{experience}

\begin{propriete}[Réaction avec la liqueur de Fehling]
L'aldéhyde réduit les ions \ce{Cu^2+} en \ce{Cu+} :
\[
\ce{R-CHO + 2 Cu^2+ + 5 HO- -> R-COO- + Cu2O v + 3 H2O}
\]
Les cétones ne réagissent pas.
\end{propriete}

\begin{remarque}
Ce test est spécifique des aldéhydes. Le précipité rouge brique est caractéristique.
\end{remarque}

\subsection{Réactif de Tollens (miroir d'argent)}

\begin{experience}[Test au réactif de Tollens]
Le réactif de Tollens est une solution ammoniacale d'oxyde d'argent \ce{[Ag(NH3)2]+}. En présence d'un aldéhyde, chauffé, il se forme un dépôt d'argent métallique sur les parois du tube à essai (miroir d'argent).
\end{experience}

\begin{propriete}[Réaction avec le réactif de Tollens]
L'aldéhyde réduit les ions \ce{Ag+} en argent métallique :
\[
\ce{R-CHO + 2 [Ag(NH3)2]+ + 2 HO- -> R-COO- + 2 Ag v + 4 NH3 + H2O}
\]
Les cétones ne réagissent pas.
\end{propriete}

\begin{aretenir}
\begin{itemize}
    \item Test à la 2,4-DNPH : positif pour tous les composés carbonylés (précipité jaune-orange).
    \item Test de Schiff, Fehling, Tollens : positifs uniquement pour les aldéhydes.
    \item Ces trois derniers tests permettent de distinguer un aldéhyde d'une cétone.
\end{itemize}
\end{aretenir}

\section{Réactions d'addition nucléophile}

\subsection{Mécanisme général}

\begin{propriete}[Addition nucléophile sur le groupe carbonyle]
Le carbone du groupe \ce{C=O}, électrophile ($\delta^+$), est attaqué par un nucléophile \ce{Nu-} (ou \ce{Nu}$:$). Il se forme un alcoolate intermédiaire, qui est ensuite protoné pour donner un alcool.
\end{propriete}

\begin{illus}
\chemfig{C(=[1]O)-[2]R_1-[0]R_2}
\hspace{1cm}
\chemfig{C(-[1]O^{-})-[2]R_1-[0]R_2}
\hspace{1cm}
\chemfig{C(-[1]OH)(-[2]R_1)-[0]R_2}
\fig{Mécanisme général de l'addition nucléophile : attaque du nucléophile \ce{Nu-} sur le carbone, formation d'un alcoolate, puis protonation.}
\end{illus}

\subsection{Addition d'eau : hydratation}

\begin{propriete}[Hydratation des aldéhydes et cétones]
L'addition d'eau sur le groupe carbonyle conduit à un gem-diol (diol vicinal sur le même carbone). La réaction est réversible.
\[
\ce{R-CO-R$'$ + H2O <=> R-C(OH)2-R$'$}
\]
L'équilibre est généralement déplacé vers le gem-diol pour les aldéhydes (surtout le méthanal), mais vers le carbonyle pour les cétones.
\end{propriete}

\subsection{Addition d'alcools : formation d'hémiacétals et d'acétals}

\begin{propriete}[Addition d'un alcool]
Un aldéhyde ou une cétone réagit avec un alcool \ce{ROH} pour former un hémiacétal (ou hémicétal), puis, en présence d'un excès d'alcool et d'un catalyseur acide, un acétal (ou cétal).
\[
\ce{R-CHO + R$'$OH <=> R-CH(OH)(OR$'$)} \quad \text{(hémiacétal)}
\]
\[
\ce{R-CH(OH)(OR$'$) + R$'$OH <=> R-CH(OR$'$)2 + H2O} \quad \text{(acétal)}
\]
\end{propriete}

\begin{remarque}
Cette réaction est importante en chimie des sucres (formes cycliques des oses).
\end{remarque}

\subsection{Addition de dérivés azotés}

\begin{propriete}[Addition de l'ammoniac et des amines]
Les aldéhydes et cétones réagissent avec l'ammoniac \ce{NH3} ou les amines primaires \ce{RNH2} pour former des imines (bases de Schiff) :
\[
\ce{R-CHO + R$'$NH2 -> R-CH=N-R$'$ + H2O}
\]
\end{propriete}

\subsection{Addition de cyanure d'hydrogène}

\begin{propriete}[Cyanhydrine]
L'addition de \ce{HCN} sur un aldéhyde ou une cétone donne une cyanhydrine (hydroxynitrile) :
\[
\ce{R-CO-R$'$ + HCN -> R-C(OH)(CN)-R$'$}
\]
Cette réaction est catalysée par une base et permet d'allonger la chaîne carbonée.
\end{propriete}

\section{Oxydation des aldéhydes}

\begin{propriete}[Oxydation des aldéhydes en acides carboxyliques]
Les aldéhydes sont facilement oxydés en acides carboxyliques par des oxydants doux comme le réactif de Tollens, la liqueur de Fehling, ou le permanganate de potassium.
\[
\ce{R-CHO + [O] -> R-COOH}
\]
\end{propriete}

\begin{remarque}
Les cétones ne sont pas oxydées par ces oxydants doux. Cette différence est utilisée pour les distinguer.
\end{remarque}

\begin{exemple}
L'éthanal \ce{CH3CHO} est oxydé en acide éthanoïque \ce{CH3COOH} par le permanganate de potassium en milieu acide :
\[
\ce{5 CH3CHO + 2 MnO4- + 6 H3O+ -> 5 CH3COOH + 2 Mn^2+ + 9 H2O}
\]
\end{exemple}

\section{Distinction aldéhyde/cétone}

\begin{aretenir}
Pour distinguer un aldéhyde d'une cétone, on utilise les tests suivants :
\begin{itemize}
    \item Test au réactif de Schiff : coloration rose pour un aldéhyde.
    \item Test à la liqueur de Fehling : précipité rouge brique pour un aldéhyde.
    \item Test au réactif de Tollens : miroir d'argent pour un aldéhyde.
    \item Les cétones ne donnent aucune de ces réactions.
\end{itemize}
\end{aretenir}

\section{L'essentiel à retenir}

\begin{synthese}
\subsection*{Définitions}
\begin{itemize}
    \item \textbf{Aldéhyde} : \ce{R-CHO} (groupe formyle en bout de chaîne).
    \item \textbf{Cétone} : \ce{R-CO-R$'$} (groupe carbonyle à l'intérieur de la chaîne).
\end{itemize}

\subsection*{Nomenclature}
\begin{itemize}
    \item Aldéhyde : suffixe \textit{-al} (ex : éthanal).
    \item Cétone : suffixe \textit{-one} avec position (ex : butan-2-one).
\end{itemize}

\subsection*{Obtention}
\begin{itemize}
    \item Alcool primaire $\xrightarrow{[O]}$ aldéhyde $\xrightarrow{[O]}$ acide carboxylique.
    \item Alcool secondaire $\xrightarrow{[O]}$ cétone.
\end{itemize}

\subsection*{Tests d'identification}
\begin{itemize}
    \item \textbf{2,4-DNPH} : précipité jaune-orange pour tout composé carbonylé.
    \item \textbf{Schiff} : coloration rose pour un aldéhyde.
    \item \textbf{Fehling} : précipité rouge brique \ce{Cu2O} pour un aldéhyde.
    \item \textbf{Tollens} : miroir d'argent pour un aldéhyde.
\end{itemize}

\subsection*{Réactions}
\begin{itemize}
    \item \textbf{Addition nucléophile} : le carbone $\delta^+$ du \ce{C=O} est attaqué par un nucléophile.
    \item \textbf{Hydratation} : formation de gem-diol (réversible).
    \item \textbf{Addition d'alcool} : hémiacétal puis acétal.
    \item \textbf{Addition de \ce{HCN}} : cyanhydrine.
    \item \textbf{Oxydation des aldéhydes} : en acide carboxylique (les cétones résistent).
\end{itemize}

\subsection*{Pièges à éviter}
\begin{itemize}
    \item Ne pas confondre aldéhyde et cétone : le test de Schiff, Fehling ou Tollens permet de les distinguer.
    \item L'oxydation d'un alcool primaire ne s'arrête pas toujours à l'aldéhyde : il faut contrôler les conditions.
    \item La 2,4-DNPH ne distingue pas les deux familles.
\end{itemize}
\end{synthese}

\section{Méthodes \& savoir-faire}

\methode{Reconnaître un aldéhyde ou une cétone à partir de sa formule}
Pour identifier un composé carbonylé, repérez le groupe \ce{C=O} :
\begin{itemize}
    \item Si le carbone du \ce{C=O} est lié à au moins un H : c'est un aldéhyde.
    \item S'il est lié à deux carbones : c'est une cétone.
\end{itemize}
\begin{exemple}
\ce{CH3CH2COCH3} : le carbone du \ce{C=O} est lié à \ce{CH3CH2-} et \ce{CH3-} (deux carbones) → c'est une cétone (pentan-2-one).
\end{exemple}

\methode{Nommer un aldéhyde ou une cétone}
\begin{enumerate}
    \item Trouver la chaîne carbonée la plus longue contenant le groupe \ce{C=O}.
    \item Pour un aldéhyde : remplacer le \textit{-e} de l'alcane par \textit{-al} (le carbone du \ce{CHO} est le n°1).
    \item Pour une cétone : remplacer le \textit{-e} par \textit{-one} et indiquer la position du \ce{C=O} par le plus petit numéro.
    \item Nommer et numéroter les substituants.
\end{enumerate}
\begin{exemple}
\chemfig{CH_3-CH_2-CH(-[:30]CH_3)-CHO} : chaîne principale = 4 carbones (butane), groupe \ce{CHO} en position 1, méthyle en position 2 → 2-méthylbutanal.
\end{exemple}

\methode{Écrire l'équation d'oxydation d'un alcool}
\begin{enumerate}
    \item Identifier le type d'alcool (primaire, secondaire, tertiaire).
    \item Pour un alcool primaire : \ce{R-CH2OH + [O] -> R-CHO + H2O}.
    \item Pour un alcool secondaire : \ce{R-CHOH-R$'$ + [O] -> R-CO-R$'$ + H2O}.
    \item Équilibrer l'équation avec l'oxydant choisi (ex : \ce{Cr2O7^2-} ou \ce{MnO4-}).
\end{enumerate}
\begin{exemple}
Oxydation du propan-1-ol \ce{CH3CH2CH2OH} en propanal \ce{CH3CH2CHO} par le dichromate :
\[
\ce{3 CH3CH2CH2OH + Cr2O7^2- + 8 H3O+ -> 3 CH3CH2CHO + 2 Cr^3+ + 15 H2O}
\]
\end{exemple}

\methode{Distinguer un aldéhyde d'une cétone par un test chimique}
\begin{enumerate}
    \item Prélever un échantillon du composé inconnu.
    \item Ajouter le réactif de Schiff (ou Fehling, ou Tollens).
    \item Chauffer si nécessaire (Fehling, Tollens).
    \item Observer : coloration rose (Schiff), précipité rouge brique (Fehling), miroir d'argent (Tollens) → aldéhyde. Aucune réaction → cétone.
\end{enumerate}
\begin{exemple}
Un composé X donne un précipité rouge brique avec la liqueur de Fehling. X est un aldéhyde.
\end{exemple}

\methode{Écrire le mécanisme d'addition nucléophile sur un carbonyle}
\begin{enumerate}
    \item Le nucléophile \ce{Nu-} attaque le carbone $\delta^+$ du \ce{C=O}.
    \item La double liaison \ce{C=O} se rompt, l'oxygène récupère une charge négative (formation d'un alcoolate).
    \item L'alcoolate est protoné par une source de protons (eau, acide) pour donner l'alcool.
\end{enumerate}
\begin{exemple}
Addition de \ce{HCN} sur l'éthanal :
\[
\ce{CH3CHO + HCN -> CH3CH(OH)(CN)}
\]
Mécanisme : \ce{CN-} attaque le carbone, puis \ce{H+} protoné l'oxygène.
\end{exemple}

\methode{Calculer la masse molaire d'un composé carbonylé}
\begin{enumerate}
    \item Écrire la formule brute du composé.
    \item Additionner les masses molaires atomiques : \ce{C} = \SI{12}{\gram\per\mol}, \ce{H} = \SI{1}{\gram\per\mol}, \ce{O} = \SI{16}{\gram\per\mol}.
\end{enumerate}
\begin{exemple}
Pour le butan-2-one \ce{C4H8O} : $M = 4 \times 12 + 8 \times 1 + 16 = \SI{72}{\gram\per\mol}$.
\end{exemple}

\methode{Prévoir le produit d'une réaction d'oxydation d'un aldéhyde}
\begin{enumerate}
    \item L'aldéhyde \ce{R-CHO} est oxydé en acide carboxylique \ce{R-COOH}.
    \item Équilibrer l'équation avec l'oxydant (ex : \ce{MnO4-} en milieu acide).
\end{enumerate}
\begin{exemple}
Oxydation du benzaldéhyde \ce{C6H5CHO} en acide benzoïque \ce{C6H5COOH} :
\[
\ce{5 C6H5CHO + 2 MnO4- + 6 H3O+ -> 5 C6H5COOH + 2 Mn^2+ + 9 H2O}
\]
\end{exemple}

\methode{Utiliser la 2,4-DNPH pour confirmer la présence d'un carbonyle}
\begin{enumerate}
    \item Préparer une solution du composé à tester.
    \item Ajouter quelques gouttes de réactif de 2,4-DNPH.
    \item Observer la formation d'un précipité jaune à orange.
\end{enumerate}
\begin{exemple}
Un composé inconnu donne un précipité orange avec la 2,4-DNPH. Il contient un groupe carbonyle (aldéhyde ou cétone).
\end{exemple}

\methode{Distinguer un aldéhyde d'une cétone par oxydation}
\begin{enumerate}
    \item Traiter le composé avec un oxydant doux (ex : \ce{KMnO4} dilué à froid).
    \item Si le composé est un aldéhyde, il est oxydé en acide carboxylique (test au papier pH : acidification).
    \item Si c'est une cétone, aucune réaction.
\end{enumerate}
\begin{exemple}
Un composé X ne réagit pas avec \ce{KMnO4} dilué à froid. X est une cétone.
\end{exemple}